% Part: first-order-logic
% Chapter: syntax-and-semantics
% Section: assignments

\documentclass[../../../include/open-logic-section]{subfiles}

\begin{document}

\olfileid{fol}{syn}{ass}

\olsection{متغیر قدر بندیاں}

\begin{explain}
!!{variable} قدر بندی~$s$ \emph{ہر} متغیر نوں قدر دیندی اے---تے
متغیر لامتناہی نیں۔ ظاہر اے کہ ایہ لازم نہیں۔ اسیں !!{variable}
قدر بندیاں توں سارے !!{variable}s نوں قدراں دیون دا تقاضا صرف
ایس لئی کردے آں کہ گل بہت سوکھی ہو جاندی اے۔ ٹرم~$t$ دی قدر
تے !!a{structure} وچ !!a{formula}~$!A$ دا~$s$ دے لحاظ نال
پورا اترنا یا نہ اترنا صرف انہاں قدراں اُتے منحصر اے جیہڑیاں
قدر بندی~$s$،~$t$ دے !!{variable}s تے~$!A$ دے آزاد
!!{variable}s نوں دیندی اے۔ اگلے دو قضیے ایہی گل دسدے نیں۔
``منحصر ہون'' نوں پکا کرن لئی اسیں ثابت کردے آں کہ دو متغیر
قدر بندیاں جے~$t$ دے سارے متغیراں اُتے متفق ہون تاں اک ہی
قدر دیندیاں نیں؛ تے جے اوہ~$!A$ دے سارے آزاد متغیراں اُتے
متفق ہون تاں $!A$ پہلی دے لحاظ نال اُدوں تے صرف اُدوں
پورا اتردا اے جدوں دوجی دے لحاظ نال پورا اتردا اے۔
\end{explain}

\begin{prop}\ollabel{prop:valindep}
جے ٹرم~$t$ دے !!{variable}s، $x_1$، \dots،~$x_n$ وچوں ہون، تے
$s_1(x_i) = s_2(x_i)$، $i = 1$، \dots،~$n$ لئی، تاں
$\Value{t}{M}[s_1] = \Value{t}{M}[s_2]$۔
\end{prop}

\begin{proof}
~$t$ دی پیچیدگی اُتے استقرا کردے آں۔ بنیادی صورت وچ $t$ یا
!!a{constant} ہو سکدا اے یا~$x_1$، \dots،~$x_n$ وچوں کوئی
متغیر۔ جے $t = c$، تاں $\Value{t}{M}[s_1] = \Assign{c}{M} =
\Value{t}{M}[s_2]$۔ جے $t = x_i$، تاں قضیے دی مفروضہ شرط
مطابق $s_1(x_i) = s_2(x_i)$، ایس لئی
$\Value{t}{M}[s_1] = s_1(x_i) = s_2(x_i) = \Value{t}{M}[s_2]$۔

استقرائی قدم لئی فرض کرو $t = \Atom{f}{t_1, \dots, t_k}$، تے
دعویٰ $t_1$، \dots، $t_k$ لئی درست اے۔ فیر
\begin{align*}
  \Value{t}{M}[s_1] & = \Value{\Atom{f}{t_1, \dots, t_k}}{M}[s_1] \\
  & = \Assign{f}{M}(\Value{t_1}{M}[s_1], \dots, \Value{t_k}{M}[s_1]).
\intertext{$j = 1$، \dots،~$k$ لئی~$t_j$ دے !!{variable}s،
    $x_1$، \dots،~$x_n$ وچوں نیں۔ استقرائی مفروضے مطابق
    $\Value{t_j}{M}[s_1] = \Value{t_j}{M}[s_2]$۔ ایس لئی}
\Value{t}{M}[s_1] & = \Value{\Atom{f}{t_1, \dots, t_k}}{M}[s_1] \\
  & = \Assign{f}{M}(\Value{t_1}{M}[s_1], \dots, \Value{t_k}{M}[s_1]) \\
  & = \Assign{f}{M}(\Value{t_1}{M}[s_2], \dots, \Value{t_k}{M}[s_2]) \\
  & = \Value{\Atom{f}{t_1, \dots, t_k}}{M}[s_2] = \Value{t}{M}[s_2].
\end{align*}
\end{proof}

\begin{prop}\ollabel{prop:satindep}
جے $!A$ دے آزاد !!{variable}s، $x_1$، \dots،~$x_n$ وچوں ہون، تے
$s_1(x_i) = s_2(x_i)$، $i = 1$، \dots،~$n$ لئی، تاں
$\Sat{M}{!A}[s_1]$ اُدوں تے صرف اُدوں جدوں $\Sat{M}{!A}[s_2]$۔
\end{prop}

\begin{proof}
اسیں $!A$ دی پیچیدگی اُتے استقرا کردے آں۔ بنیادی صورت وچ، جدوں
$!A$ ایٹمی ہووے، $!A$ ایہناں وچوں ہو سکدا اے:
\iftag{prvTrue}{$\ltrue$،}{}
\iftag{prvFalse}{$\lfalse$،}{}
$\Atom{R}{t_1, \dots, t_k}$، جیہڑا $k$-مقامی محمول $R$ تے
ٹرماں $t_1$، \dots،~$t_k$ نال بندا اے؛ یا $\eq[t_1][t_2]$،
جیہڑا ٹرماں $t_1$ تے~$t_2$ نال بندا اے۔ پچھلی دو صورتاں وچ اسیں !!{biconditional}
دی صرف اگلی سمت دکھاندے آں، کیوں جے اُلٹی سمت دا ثبوت
متناظر اے۔

\begin{enumerate}
\tagitem{prvTrue}{%
  \indcase{!A}{\ltrue}{$\Sat{M}{\indfrm}[s_1]$ تے
    $\Sat{M}{\indfrm}[s_2]$ دونیں۔}}{}

\tagitem{prvFalse}{%
  \indcase{!A}{\lfalse}{$\Sat/{M}{!A}[s_1]$ تے
    $\Sat/{M}{!A}[s_2]$ دونیں۔}}{}

\item
  \indcase{!A}{\Atom{R}{t_1, \ldots, t_k}}{فرض کرو
    $\Sat{M}{\indfrm}[s_1]$۔ تاں
    \[
    \langle \Value{t_1}{M}[s_1], \ldots, \Value{t_k}{M}[s_1] \rangle
    \in \Assign{R}{M}.
    \]
    $i = 1$، \dots،~$k$ لئی \olref{prop:valindep} مطابق
    $\Value{t_i}{M}[s_1] = \Value{t_i}{M}[s_2]$۔ ایس لئی
    % PNB-SEM-011: ٹپل دی پہلی جگہ پہلا ٹرم ای رہندا اے، چلدا اشاریہ نہیں۔
    $\langle \Value{t_1}{M}[s_2], \ldots, \Value{t_k}{M}[s_2] \rangle
    \in \Assign{R}{M}$، تے نتیجتاً $\Sat{M}{\indfrm}[s_2]$۔}
\item
  \indcase{!A}{\eq[t_1][t_2]}{فرض کرو $\Sat{M}{\indfrm}[s_1]$۔
    تاں $\Value{t_1}{M}[s_1] = \Value{t_2}{M}[s_1]$۔ ایس لئی
    \begin{align*}
      \Value{t_1}{M}[s_2] & = \Value{t_1}{M}[s_1]
      & \text{(\olref{prop:valindep} مطابق)} \\
      & = \Value{t_2}{M}[s_1]
      & \text{(کیوں جے $\Sat{M}{\eq[t_1][t_2]}[s_1]$)}\\
      &= \Value{t_2}{M}[s_2]
      & \text{(\olref{prop:valindep} مطابق)،}
    \end{align*}
    ایس لئی $\Sat{M}{\eq[t_1][t_2]}[s_2]$۔}
\end{enumerate}

ہُن فرض کرو کہ $\Sat{M}{!B}[s_1]$ اُدوں تے صرف اُدوں جدوں
$\Sat{M}{!B}[s_2]$، سارے اوہناں !!{formula}s~$!B$ لئی
جیہڑے~$!A$ نالوں گھٹ پیچیدہ نیں۔ استقرائی قدم~$!A$ دے مرکزی عامل مطابق
صورتاں وچ ونڈیا جاندا اے۔ ہر صورت وچ اسیں !!{biconditional}
دی صرف اگلی سمت دکھاندے آں؛ اُلٹی سمت دا ثبوت متناظر اے۔
مسوّراں والیاں صورتاں توں چھٹ اسیں استقرائی مفروضہ ذیلی-!!{formula}s~$!B$
اُتے لاگو کردے آں، جیہڑے~$!A$ دے نیں۔~$!B$ دے آزاد متغیر،
~$!A$ دے آزاد متغیراں وچوں نیں۔ ایس لئی جے $s_1$ تے $s_2$،~$!A$
دے آزاد متغیراں اُتے متفق ہون، تاں اوہ~$!B$ دے آزاد متغیراں
اُتے وی متفق ہون گے، تے استقرائی مفروضہ~$!B$ اُتے لاگو ہووے گا۔

\begin{enumerate}
\tagitem{defNot}{}{%
  \iftag{probNot}{%
    \indcase!{!A}{\lnot !B}{}}{%
    \indcase{!A}{\lnot !B}{جے $\Sat{M}{\indfrm}[s_1]$، تاں
      $\Sat/{M}{!B}[s_1]$؛ استقرائی مفروضے مطابق
      $\Sat/{M}{!B}[s_2]$، لہٰذا $\Sat{M}{\indfrm}[s_2]$۔}}}

\tagitem{defAnd}{}{%
  \iftag{probAnd}{%
    \indcase!{!A}{!B \land !C}{}}{%
    \indcase{!A}{!B \land !C}{جے $\Sat{M}{\indfrm}[s_1]$، تاں
      $\Sat{M}{!B}[s_1]$ تے $\Sat{M}{!C}[s_1]$؛ استقرائی
      مفروضے مطابق $\Sat{M}{!B}[s_2]$ تے $\Sat{M}{!C}[s_2]$۔ ایس لئی
      $\Sat{M}{\indfrm}[s_2]$۔}}}

\tagitem{defOr}{}{%
  \iftag{probOr}{%
    \indcase!{!A}{!B \lor !C}{}}{%
    \indcase{!A}{!B \lor !C}{جے $\Sat{M}{\indfrm}[s_1]$، تاں
      $\Sat{M}{!B}[s_1]$ یا $\Sat{M}{!C}[s_1]$۔ استقرائی مفروضے
      مطابق $\Sat{M}{!B}[s_2]$ یا $\Sat{M}{!C}[s_2]$؛ لہٰذا
      $\Sat{M}{\indfrm}[s_2]$۔}}}

\tagitem{defIf}{}{%
  \iftag{probIf}{%
    \indcase!{!A}{!B \lif !C}{}}{%
    \indcase{!A}{!B \lif !C}{جے $\Sat{M}{\indfrm}[s_1]$، تاں
    $\Sat/{M}{!B}[s_1]$ یا $\Sat{M}{!C}[s_1]$۔ استقرائی مفروضے
    مطابق $\Sat/{M}{!B}[s_2]$ یا $\Sat{M}{!C}[s_2]$؛ ایس لئی
    $\Sat{M}{\indfrm}[s_2]$۔}}}

\tagitem{defIff}{}{%
  \iftag{probIff}{%
    \indcase!{!A}{!B \liff !C}{}}{%
    \indcase{!A}{!B \liff !C}{جے $\Sat{M}{\indfrm}[s_1]$، تاں یا تاں
      $\Sat{M}{!B}[s_1]$ تے $\Sat{M}{!C}[s_1]$ دونیں، یا $\Sat/{M}{!B}[s_1]$ تے
      $\Sat/{M}{!C}[s_1]$ دونیں۔ استقرائی مفروضے مطابق یا تاں
      $\Sat{M}{!B}[s_2]$ تے $\Sat{M}{!C}[s_2]$ دونیں، یا $\Sat/{M}{!B}[s_2]$
      تے $\Sat/{M}{!C}[s_2]$ دونیں۔ دونوں صورتاں وچ $\Sat{M}{\indfrm}[s_2]$۔}}}

\tagitem{defEx}{}{%
  \iftag{probEx}{%
    \indcase!{!A}{\lexists[x][!B]}{}}{%
    \indcase{!A}{\lexists[x][!B]}{جے $\Sat{M}{\indfrm}[s_1]$، تاں
      کوئی $m \in \Domain{M}$ ایسا اے جیہدے لئی
      $\Sat{M}{!B}[\Subst{s_1}{m}{x}]$۔ رکھو
      $s_1' = \Subst{s_1}{m}{x}$ تے $s_2' =
      \Subst{s_2}{m}{x}$۔~$!B$ دے آزاد متغیر $x_1$،
      \dots، $x_n$ تے $x$ وچوں نیں۔ $s_1'(x_i) = s_2'(x_i)$، کیوں جے
      $s_1'$ تے $s_2'$، $x$-بدل قدر بندیاں نیں، بالترتیب $s_1$ تے~$s_2$ دیاں،
      تے مفروضے مطابق $s_1(x_i) = s_2(x_i)$۔
      % PNB-SEM-018: درج متغیر مسوّر دے متغیر دے برابر وی ہو سکدا اے؛ دونوں صورتاں دی برابری صریح اے۔
      ایہ سبب مسوّر دے متغیر توں ہور درج متغیراں لئی اے؛ جے کوئی
      درج متغیر مسوّر دے متغیر دے برابر ہووے، تاں دوہاں وچ رکھی
      اکّو قدر نال برابری ملدی اے۔ انہاں دی تعریف توں
      $s_1'(x) = s_2'(x) = m$؛ $s_1'$ تے~$s_2'$ دی تعریف وچ ایہی
      قدر رکھی گئی سی۔ فیر استقرائی
      مفروضہ $!B$ تے $s_1'$، $s_2'$ اُتے لاگو ہوندا اے، ایس لئی
      $\Sat{M}{!B}[s_2']$۔ ہُن چونکہ $s_2' = \Subst{s_2}{m}{x}$،
      کوئی $m \in \Domain{M}$ ایسا اے جیہدے لئی
      $\Sat{M}{!B}[\Subst{s_2}{m}{x}]$، لہٰذا
      $\Sat{M}{\indfrm}[s_2]$۔}}}

\tagitem{defAll}{}{%
  \iftag{probAll}{%
    \indcase!{!A}{\lforall[x][!B]}{}}{%
    \indcase{!A}{\lforall[x][!B]}{جے $\Sat{M}{\indfrm}[s_1]$، تاں ہر
      $m \in \Domain{M}$ لئی $\Sat{M}{!B}[\Subst{s_1}{m}{x}]$۔ اسیں
      دکھانا چاہندے آں کہ ہر $m \in \Domain{M}$ لئی
      $\Sat{M}{!B}[\Subst{s_2}{m}{x}]$ وی اے۔ سو کوئی وی
      $m \in \Domain{M}$ چنو، تے
      % PNB-SEM-012: ماخذ وچ اک غیر متعین قدر بندی دو وار بدل دی اے؛ دونوں بدل صورتاں اپنی اپنی اصل قدر بندی توں بندیاں نیں۔
      $s_1' = \Subst{s_1}{m}{x}$ تے $s_2' =
      \Subst{s_2}{m}{x}$ ویکھو۔ ساڈے کول $\Sat{M}{!B}[s_1']$ اے۔
      ~ $!B$ دے آزاد متغیر $x_1$، \dots، $x_n$ تے $x$ وچوں نیں۔
      $s_1'(x_i) = s_2'(x_i)$، کیوں جے $s_1'$ تے $s_2'$
      $x$-بدل قدر بندیاں نیں، بالترتیب $s_1$ تے~$s_2$ دیاں، تے مفروضے مطابق
      $s_1(x_i) = s_2(x_i)$۔ ایہ سبب مسوّر دے متغیر توں ہور
      درج متغیراں لئی اے؛ جے کوئی درج متغیر مسوّر دے متغیر دے
      برابر ہووے، تاں دوہاں وچ رکھی اکّو قدر نال برابری ملدی اے۔
      تعریف توں $s_1'(x) = s_2'(x) = m$؛
      ایہ~$s_1'$ تے~$s_2'$ نوں دتی گئی قدر اے۔ فیر استقرائی مفروضہ~$!B$
      تے~$s_1'$، $s_2'$ اُتے لاگو ہوندا اے، تے
      $\Sat{M}{!B}[s_2']$ ملدا اے۔ ایہ ہر $m \in \Domain{M}$ لئی اے؛
      یعنی $\Sat{M}{!B}[\Subst{s_2}{m}{x}]$ سارے~$m \in
      \Domain{M}$ لئی اے۔ ایس لئی $\Sat{M}{\indfrm}[s_2]$۔}}}
\end{enumerate}
استقرا نال نتیجہ ملدا اے کہ $\Sat{M}{!A}[s_1]$ اُدوں تے صرف اُدوں
جدوں $\Sat{M}{!A}[s_2]$، جے $!A$ دے آزاد !!{variable}s، $x_1$،
\dots، $x_n$ وچوں ہون تے $s_1(x_i)=s_2(x_i)$،
$i = 1$، \dots،~$n$ لئی۔
\end{proof}

\begin{probtag}{probNot,probOr,probAnd,probIf,probIff,probEx,probAll}
\olref[fol][syn][ass]{prop:satindep} دا ثبوت مکمل کرو۔
\end{probtag}

\begin{explain}
!!^{sentence}s وچ آزاد متغیر نہیں ہوندے؛ ایس لئی کوئی وی دو متغیر
قدر بندیاں کسے وی جملے دے سارے آزاد متغیراں (جنہاں دی گنتی صفر
اے) نوں اوہی قدراں دیندیاں نیں۔ پچھلا ثابت کیتا قضیہ ایہ دسدا اے
کہ !!a{sentence} دا کسے ساخت وچ متغیر قدر بندی دے لحاظ نال
پورا اترنا یا نہ اترنا اُس قدر بندی توں بالکل آزاد اے۔ اسیں ایہ
حقیقت درج کردے آں۔ ایہی اگلی تعریف نوں جائز بناندی اے، جیہڑی
!!a{structure} وچ !!a{sentence} دے پورا اترن لئی متغیر قدر
بندی دا ذکر نہیں کردی۔
\end{explain}

\begin{cor}
\ollabel{cor:sat-sentence}
جے $!A$ اک !!a{sentence} اے تے $s$ متغیر قدر بندی اے، تاں
$\Sat{M}{!A}[s]$ اُدوں تے صرف اُدوں جدوں
$\Sat{M}{!A}[s']$، ہر متغیر قدر بندی~$s'$ لئی۔
\end{cor}

\begin{proof}
  $s'$ کوئی وی متغیر قدر بندی لو۔ چونکہ $!A$ اک !!a{sentence} اے،
  اوہدے کوئی آزاد متغیر نہیں؛ ایس لئی ہر متغیر قدر بندی~$s'$
  بدیہی طور تے~$!A$ دے سارے آزاد متغیراں نوں اوہی قدراں دیندی
  اے جیہڑیاں~$s$ دیندی اے۔ سو \olref{prop:satindep} دی شرط پوری
  اے، تے $\Sat{M}{!A}[s]$ اُدوں تے صرف اُدوں جدوں
  $\Sat{M}{!A}[s']$۔
\end{proof}

\begin{defn}
\ollabel{defn:satisfaction}
جے $!A$ اک !!a{sentence} اے، تاں اسیں آکھدے آں کہ
!!a{structure}~$\Struct M$،~$!A$ نوں \emph{پورا کردی اے}،
$\Sat{M}{!A}$، اُدوں تے صرف اُدوں جدوں
$\Sat{M}{!A}[s]$ ہر متغیر قدر بندی~$s$ لئی۔
\end{defn}

جے $\Sat{M}{!A}$، تاں اسیں سادہ لفظاں وچ ایہ وی آکھدے آں کہ
\emph{$!A$،~$\Struct{M}$ وچ سچ اے۔} پورا اترن دا تصور
انفرادی !!{sentence}s توں !!{sentence}s دے سیٹاں تک فطری طور تے
پھیل جاندا اے۔

\begin{defn}
\ollabel{defn:sat}
جے $\Gamma$، !!{sentence}s دا سیٹ~$\Gamma$ اے، تاں اسیں
آکھدے آں کہ !!a{structure}~$\Struct M$،~$\Gamma$ نوں
\emph{پورا کردی اے}، $\Sat{M}{\Gamma}$، اُدوں تے صرف اُدوں
جدوں $\Sat{M}{!A}$ ہر $!A \in \Gamma$ لئی۔
\end{defn}

\begin{prop}\ollabel{prop:sentence-sat-true}
  فرض کرو $\Struct{M}$ اک !!a{structure}، $!A$ اک
  !!a{sentence} تے $s$ متغیر قدر بندی اے۔ $\Sat{M}{!A}$
  اُدوں تے صرف اُدوں جدوں $\Sat{M}{!A}[s]$۔
\end{prop}

\begin{proof}
مشق۔
\end{proof}

\begin{prob}
\olref[fol][syn][ass]{prop:sentence-sat-true} ثابت کرو۔
\end{prob}

\begin{prop}\ollabel{prop:sat-quant}
  فرض کرو $!A(x)$ وچ صرف $x$ آزاد اے، تے $\Struct M$ اک
  !!a{structure} اے۔ فیر:
  \begin{tagenumerate}{prvEx,prvAll}
    \tagitem{prvEx}{$\Sat{M}{\lexists[x][!A(x)]}$ اُدوں تے صرف اُدوں جدوں
      $\Sat{M}{!A(x)}[s]$ گھٹ توں گھٹ اک متغیر قدر بندی~$s$ لئی۔}{}
    \tagitem{prvAll}{$\Sat{M}{\lforall[x][!A(x)]}$ اُدوں تے صرف اُدوں جدوں
      $\Sat{M}{!A(x)}[s]$ ہر متغیر قدر بندی~$s$ لئی۔}{}
\end{tagenumerate}
\end{prop}

\begin{proof}
مشق۔
\end{proof}

\begin{prob}
\olref[fol][syn][ass]{prop:sat-quant} ثابت کرو۔
\end{prob}

\begin{prob}
\DeclareRobustCommand{\VDash}{\mathrel{||}\joinrel\Relbar}
فرض کرو $\Lang L$ ایسی زبان اے جیہدے وچ !!{function}s نہیں۔
% PNB-SEM-021: تازہ ثابتہ نشان چنن لئی معیاری زبان دا لامتناہی ذخیرہ ایتھے صریح مفروضہ اے۔
ایہ مشق ثابتہ نشاناں دا لامتناہی ذخیرہ من کے کیتی جاندی اے،
تاں جو ہر قدم اُتے تازہ ناں چنیا جا سکے۔
!!{structure}~$\Struct M$، !!a{constant}~$c$ تے
$a \in \Domain M$ دتے ہون، تاں $\Struct M[a/c]$ اوہ
!!{structure} تعریف کرو جیہڑی~$\Struct M$ ورگی اے، سوائے ایس دے
کہ $\Assign{c}{M[a/c]} = a$۔ $\Struct M \VDash !A$،
!!{sentence}s~$!A$ لئی ایویں تعریف کرو:
\begin{enumerate}
\tagitem{prvFalse}{%
  \indcase{!A}{\lfalse}{$\Struct{M} \VDash \indfrm$ نہیں۔}}{}

\tagitem{prvTrue}{%
  \indcase{!A}{\ltrue}{$\Struct{M} \VDash \indfrm$۔}}{}

\item \indcase{!A}{\Atom{R}{d_1, \dots, d_n}}{$\Struct M \VDash \indfrm$
  اُدوں تے صرف اُدوں جدوں $\langle \Assign{d_1}{M}, \dots, \Assign{d_n}{M} \rangle \in
  \Assign{R}{M}$۔}

\item \indcase{!A}{\eq[d_1][d_2]}{$\Struct M \VDash \indfrm$ اُدوں تے صرف اُدوں جدوں
  $\Assign{d_1}{M} = \Assign{d_2}{M}$۔}

\tagitem{prvNot}{%
  \indcase{!A}{\lnot !B}{$\Struct{M} \VDash \indfrm$ اُدوں تے صرف اُدوں جدوں
    $\Struct{M} \VDash {!B}$ نہ ہووے۔}}{}

\tagitem{prvAnd}{%
  \indcase{!A}{(!B \land !C)}{$\Struct{M} \VDash
    \indfrm$ اُدوں تے صرف اُدوں جدوں $\Struct{M} \VDash!B$ تے
    $\Struct{M} \VDash !C$ دونیں۔}}{}

\tagitem{prvOr}{%
  \indcase{!A}{(!B \lor !C)}{$\Struct{M} \VDash \indfrm$ اُدوں تے صرف اُدوں جدوں
    $\Struct{M} \VDash !B$ یا $\Struct{M} \VDash !C$، یا دونیں۔}}{}

\tagitem{prvIf}{%
  \indcase{!A}{(!B \lif !C)}{$\Struct{M} \VDash \indfrm$ اُدوں تے صرف اُدوں جدوں
    $\Struct {M} \VDash !B$ نہ ہووے یا $\Struct M \VDash !C$ ہووے، یا دونیں۔}}{}

\tagitem{prvIff}{%
  \indcase{!A}{(!B \liff !C)}{$\Struct{M} \VDash {\indfrm}$ اُدوں تے صرف اُدوں جدوں یا تاں
    $\Struct{M} \VDash {!B}$ تے $\Struct{M} \VDash {!C}$ دونیں، یا
    نہ $\Struct{M} \VDash {!B}$ ہووے نہ $\Struct{M} \VDash {!C}$۔}}{}

\tagitem{prvAll}{%
  \indcase{!A}{\lforall[x][!B]}{$\Struct{M} \VDash {\indfrm}$ اُدوں تے صرف اُدوں جدوں
    ہر $a \in \Domain{M}$ لئی $\Struct{M[a/c]} \VDash \Subst{!B}{c}{x}$،
    جے $c$،~$!B$ وچ نہ آوے۔}}{}

\tagitem{prvEx}{%
  \indcase{!A}{\lexists[x][!B]}{$\Struct{M} \VDash
  {\indfrm}$ اُدوں تے صرف اُدوں جدوں کوئی $a \in \Domain M$ ایسا ہووے
  جیہدے لئی $\Struct{M[a/c]} \VDash \Subst{!B}{c}{x}$، جے
  $c$،~$!B$ وچ نہ آوے۔}}{}
\end{enumerate}
$x_1$، \dots، $x_n$،~$!A$ دے سارے آزاد !!{variable}s ہون؛
$c_1$، \dots، $c_n$،~$!A$ وچ نہ آون والے ثابتہ نشان ہون؛
انہاں تازہ ثابتہ نشاناں نوں آپس وچ وکھرا چنو؛
$a_1$، \dots، $a_n \in \Domain M$ ہون؛ تے $s(x_i) = a_i$۔

ثابت کرو کہ $\Sat{M}{!A}[s]$ اُدوں تے صرف اُدوں جدوں
$\Struct M[a_1/c_1,\dots,a_n/c_n]
\VDash \Subst{\Subst{!A}{c_1}{x_1}\dots}{c_n}{x_n}$۔

(ایہ مسئلہ دکھاندا اے کہ پہلے درجے دی منطق لئی متغیر قدر بندیاں
توں بغیر وی معنویات دتی جا سکدی اے۔)
\end{prob}

\begin{prob}
فرض کرو $f$ اک یک-مقامی فنکشن نشان اے جیہڑا~$!A(x,y)$ وچ نہیں آندا۔
% PNB-SEM-022: قدر بندی توں بغیر صدق تے سکولم بدل لئی جملہ ہون تے جائز بدل دیاں پچھلی شرطاں صریح نیں۔
ایتھے فارمولے دے آزاد متغیر صرف دتے دوناں متغیراں وچوں نیں،
تے دتے فنکشن ٹرم دا بدل پچھلی تعریف مطابق بدل لئی آزاد اے۔
ثابت کرو کہ کوئی !!a{structure}~$\Struct{M}$ ایسی موجود اے جیہدے لئی
$\Sat{M}{\lforall[x][\lexists[y][!A(x,y)]]}$، اُدوں تے صرف اُدوں
جدوں کوئی~$\Struct M'$ ایسی موجود ہووے جیہدے لئی
$\Sat{M'}{\lforall[x][!A(x,f(x))]}$۔

(ایہ مسئلہ سکولم دے قضیے دی اک خاص صورت اے؛
$\lforall[x][!A(x,f(x))]$ نوں
$\lforall[x][\lexists[y][!A(x,y)]]$ دی \emph{سکولم معیاری صورت}
آکھدے آں۔)
\end{prob}

\end{document}
